\runninghead{Scope and References}
\section{Appendix: Scope and Qualifications}

\subsection{Liturgical Identity, Rights, and Study English}
This concise study treats the First Sunday of Advent, Year~B (\texttt{PC-S01-B}; \work{Ordo lectionum Missae} n.~2; U.S.\ Lectionary No.~2), opening Liturgical Year~2027 in the Dioceses of the United States of America on Sunday, November~29, 2026 (Cycle~B, Weekday~I), under the \work{Missale Romanum} (2002/2008), \work{The Roman Missal} (U.S.\ English, 2011, with the February 2021 USCCB Collect conclusion emendation), and the \work{Lectionary for Mass} (1998/2002). The seven Missal elements of \work{Dominica I Adventus} are shared across Years~A, B, and~C (owned canonically by the shared \texttt{PC-S01} Missal record), while the five Liturgy of the Word selections belong to Year~B; joining checked patristic verse-level commentaries across all twelve parts is an editorial synthesis. Full appointed study texts, branch tables, and chronological evidence notes appear in the companion expansive study.

Approved proclaimed U.S.\ liturgical texts --- \work{The Roman Missal} (English copyright ICEL, 2010/2011), protected modern Latin Missal oration, Preface, and solemn-blessing bodies, and the U.S.\ \work{Lectionary for Mass} (NABRE and Revised Grail / Abbey Psalms) --- are neither reproduced nor closely paraphrased. All English Scripture quotations come from public-domain Douay--Rheims (Challoner) strictly as \textbf{public-domain study English}, never as the approved proclaimed English of the U.S.\ Lectionary. Latin Scripture follows the Clementine Vulgate (with the Communion Antiphon \emph{Dominus dabit benignitatem, et terra nostra dabit fructum suum} also collated against the secondary 2002 digital Missal reproduction, p.~83). Hebrew/Lectionary Psalms~25, 80, and~85 equal Vulgate Psalms~24, 79, and~84; Hebrew/Lectionary Isa.~63:19b and~64:2b--7 equal Vulgate Isa.~64:1 and~64:3b--8. \emph{Scriptural Date and Location} follows the Vulgate numbering of the inspected Latin controls.

\subsection{Verification Bounds and Source-Copy Qualifications}
Primary research was conducted offline against held records without reopening restricted remote payloads. Accordingly: (1) printed 2002 and printed 2008 \work{Missale Romanum} altar-book correspondence (including Advent entries of \work{Notitiae}~44 [2008]) remains unverified offline, so comparisons for the Communion Antiphon and \work{Veronense} XXVIIII.372 (matching the secondary 2002 digital Missal copy except for one inflection of ``temporal,'' ancient \emph{temporalis} against secondary-2002 \emph{temporali}) rest on the secondary 2002 digital reproduction; (2) exact Entrance Antiphon cue/truncation, modern Missal orations, Advent Preface subtitle and date-range rubric, optional Advent solemn blessings, exact U.S.\ Lectionary No.~174 Common Psalm refrains/spans, Appendix Eucharistic Prayers, and unextracted GIRM articles remain unverified offline; (3) ancient sacramentaries (\work{Gelasianum Vetus} II.lxxxiv, Postcommunion prayer~683, and \work{Veronense} XXVIIII.372) serve strictly as comparative antecedents, never substitutes for modern prayers, and exact lexical alignment between \work{Gelasianum Vetus} II.lxxxiv and the modern 2002/2008 Collect was not collated offline (while \emph{Prosint nobis} has no established ancestor in the held ancient sacramentaries); (4) Bede (\work{In Marci Evangelium Expositio} IV.13) follows the Wikisource/Corpus Corporum transcription without PL~92 print-facsimile collation; (5) Jerome (PL~24 cols.~642--651, excluding Vallarsi's note on Origen), Aquinas (\work{Super Isaiam} Parma Vol.~14, pp.~370--371), Lapide (Vol.~7 on 1~Cor.~1:3--4, 6, 8--9; vv.~5 and~7 facsimile-checked), and \work{Gelasianum Vetus} II.lxxxiv (Wilson 1894, p.~218, preserving one unreadable OCR word before \emph{voluntatem}) rest on uncollated raw OCR copies; and (6) Augustine (\work{Ep.}~199 §3), Theophylact of Ohrid, and Pseudo-Jerome on Mark~13:32--37 are cited strictly as mediated through Aquinas's \work{Catena Aurea in Marcum} (1842 Oxford English edition).

\section*{References}

\begin{itemize}[leftmargin=1.5em,itemsep=0.25em]

\item \textbf{Biblical Texts (Study Controls).} \work{Biblia Sacra juxta Vulgatam Clementinam} (eBible \texttt{latvuc}), Isa.~63:16--17, 64:1, 3--8; Ps.~24:1--3; Ps.~79:2--4, 15--16, 18--19; Ps.~84:8, 13; 1~Cor.~1:3--9; Mark~13:32--37. \work{The Holy Bible, Douay--Rheims Version} (Challoner, Project Gutenberg~1581), public-domain study English for the same passages (not the approved U.S.\ Lectionary).

\item \textbf{Liturgical Books and Rubrics.} \work{Missale Romanum} (2002/2008), \work{Dominica I Adventus} and \work{Ordo Missae}; \work{The Roman Missal} (U.S.\ English, 2011, emended February 2021); protected modern Latin and ICEL English texts are not reproduced. \work{General Instruction of the Roman Missal} (U.S.\ English, 2011--2021), nn.~48, 53, 61, 62, 64, 68, 74, 87, 90, 364, 365a--d (summarized only). \work{Ordo lectionum Missae} (1981), n.~2; U.S.\ \work{Lectionary for Mass} (1998/2002), Nos.~2 and~174 (protected texts not reproduced).

\item \textbf{Ancient Roman Sacramentaries.} \work{Sacramentarium Gelasianum Vetus}, ed.\ H.~A.\ Wilson (Oxford, 1894), Book~II, lxxxiv, \work{Postcommunio}~683, p.~218 (uncollated raw OCR copy; exact lexical alignment with the modern 2002/2008 Collect not collated offline). \work{Sacramentarium Veronense}, ed.\ C.~L.\ Feltoe (Cambridge, 1896), Section~XXVIIII, marginal~372, p.~74 (bounded comparison against secondary 2002 digital Missal reproduction).

\item \textbf{St.~Augustine of Hippo.} \work{Enarrationes in Psalmos}, NPNF~1, Vol.~8 (1888): Ps.~24 (XXV), §§1--3, 10; Ps.~79 (LXXX), §§1--3, 9, 11; Ps.~84 (LXXXV), §§6, 10--12. \work{De gratia et libero arbitrio}, ch.~45 [XXIII], NPNF~1, Vol.~5. \work{Epistula}~199, §3, mediated via Aquinas, \work{Catena Aurea in Marcum} 13:32--37.

\item \textbf{St.~Bede the Venerable.} \work{In Marci Evangelium Expositio}, Book~IV, ch.~13 (on Mark~13:32--37), Wikisource/Corpus Corporum transcription (not print-collated against PL~92).

\item \textbf{St.~Bernard of Clairvaux.} \work{Sermones in Adventu Domini}, Sermo~3.4 and Sermo~5.1 (PL~183; 1921 English translation).

\item \textbf{Cornelius a Lapide.} \work{Commentaria in omnes divi Pauli epistolas}, Vol.~7 (1908 English), on 1~Cor.~1:3--9 (vv.~5, 7 facsimile-checked; vv.~3--4, 6, 8--9 raw OCR; citing Trent, Sess.~VI, ch.~11 and~13).

\item \textbf{St.~Cyril of Jerusalem.} \work{Catechesis}~15, §§1--2, 4, 26, NPNF~2, Vol.~7.

\item \textbf{St.~Gregory the Great.} \work{Homiliae in Evangelia}, Homily~1 (Luke~21:25--33, §§1, 5, PL~76), Homily~9 (Matt.~25:14--30), Homily~13 (Luke~12:35--40), Wikisource Latin extraction.

\item \textbf{St.~Hilary of Poitiers.} \work{De Trinitate}, Book~IX, §§58--68 (esp.\ §§65, 67, on Mark~13:32 / Matt.~24:36 with Matt.~25:13; 24:44, 46), NPNF~2, Vol.~9 (1898).

\item \textbf{St.~Jerome.} \work{Commentarii in Isaiam Prophetam}, Book~XVII (Isa.~63:16--17; 64:1, 3--8), PL~24 (1865), cols.~642--651 (raw OCR; noting St.~Irenaeus, \work{Adversus Haereses} Book~I).

\item \textbf{St.~John Chrysostom.} \work{Homilies on First Corinthians}, Homily~II (§§1--8, on 1~Cor.~1:4--9), NPNF~1, Vol.~12.

\item \textbf{St.~Leo the Great.} \work{Sermon}~19, §§1--3, NPNF~2, Vol.~12.

\item \textbf{St.~Robert Bellarmine.} \work{Explanatio in Psalmos} (trans.\ John O'Sullivan): Ps.~24, explanations~1--3; Ps.~79, explanations~1--3, 14--18; Ps.~84, explanations~7, 12.

\item \textbf{St.~Thomas Aquinas.} \work{Expositio super Isaiam ad litteram}, ch.~63--64, Parma Vol.~14 (1863), pp.~370--371 (raw OCR). \work{Catena Aurea in Marcum}, 13:32--37 (Oxford, 1842; also mediating Theophylact of Ohrid, Pseudo-Jerome, and St.~Augustine, \work{Ep.}~199 §3).

\end{itemize}
